\section{Prediction Equation}

Let's express the relationship between the vectors

\begin{equation}
	\mathbf{u} = \left[
		\begin{matrix}
			u(0) \\ 
			u(1) \\ 
			\vdots \\ 
			u(N-1)
		\end{matrix}
	\right],
	\mathbf{x} = \left[
		\begin{matrix}
			x(1) \\ 
			x(2) \\ 
			\vdots \\ 
			x(N)
		\end{matrix}
	\right]
\end{equation}

\noindent
based on the state equation

\begin{equation}
	x(k+1) = Ax(k) + Bu(k)
\end{equation}

\noindent
and initial condition $x(0)$.

The solution for the state equation is:

\begin{equation}
	x(k) = A^kx(0) + \sum_{i=0}^{k-1} A^{k-1-i} B u(i)
\end{equation}

\noindent
and it can be written for $k = 1,2,...,N$ as

\noindent
\begin{align}
	x(1) &= A x(0) + Bu(0) \\
	x(2) &= A^2 x(0) + ABu(0) + Bu(1) \\
	&\vdots \\
	x(N) &= A^N x(0) A^{N - 1}Bu(0) + \dotsi + ABu(N - 2) + Bu(N - 1)
\end{align}

\noindent
or

\begin{equation}
	\left[
		\begin{matrix}
			x(1) \\
			x(2) \\
			\vdots \\
			x(N) \\
		\end{matrix}
	\right] = 
	\left[
		\begin{matrix}
			A \\
			A^2 \\
			\vdots \\
			A^N
		\end{matrix}
	\right]
	x(0) + 
	\left[
		\begin{matrix}
			B & 0 & \dotsi & 0 \\
			AB & B & \dotsi & 0 \\
			\vdots & \vdots & \ddots & \vdots \\
			A^{N-1}B & A^{N-2}B & \dotsi & B
		\end{matrix}
	\right]
	\left[
		\begin{matrix}
			u(0) \\
			u(1)\\
			\vdots \\
			u(N-1)
		\end{matrix}
	\right]
\end{equation}

This can be written in a shorter form as

\begin{equation}
	\mathbf{x} = \mathbf{A}x(0) + \mathbf{Bu}
\end{equation}

\noindent
with 

\begin{equation}
	\mathbf{A} = \left[
		\begin{matrix}
			A \\
			A^2 \\
			\vdots \\
			A^N
		\end{matrix}
	\right]
	, 
	\mathbf{B} = \left[
		\begin{matrix}
			B & 0 & \dotsi & 0 \\
			AB & B & \dotsi & 0 \\
			\vdots & \vdots & \ddots & \vdots \\
			A^{N-1}B & A^{N-2}B & \dotsi & B
		\end{matrix}
	\right]
\end{equation}

\purple{
	
================================

// NOTE: mudar funcao custo para `$\bar{x}$ = 0` 

================================

}

Tendo em vista o v�nculo

\begin{equation}
	\mathbf{x} = \mathbf{A}x(0) + \mathbf{Bu}
\end{equation}

pode-se reescrever a fun��o de custo como

\begin{equation}
		J(\mathbf{x}, \mathbf{u}) = \norma{x(0) - \bar{x}}{Q} 
		+ \green{\overbrace{\norma{\mathbf{x} - col_N(\bar{x})}{\mathbf{Q}}}^\square} 
		+ \blue{\overbrace{\norma{\mathbf{u} - col_N(\bar{u})}{\mathbf{R}}}^\bigcirc}
\end{equation}

\noindent
\begin{align}
	\green{\square} &= \green{\norma{\mathbf{x} - col_N(\bar{x})}{\mathbf{Q}}} \\
	&= \green{
		\Big[
			\mathbf{A}x(0) + \mathbf{Bu} - col_N(\bar{x})
		\Big]^T \mathbf{Q}
		\Big[
			\mathbf{A}x(0) + \mathbf{Bu} - col_N(\bar{x})
		\Big]
	} \\
	&= \green{
		\Big[
			\mathbf{A}x(0) - col_N(\bar{x}) + \mathbf{Bu}
		\Big]^T \mathbf{Q}
		\Big[
			\mathbf{A}x(0) - col_N(\bar{x}) + \mathbf{Bu}
		\Big]
	} \\
	\begin{split}
		&= \green{
			\norma{\mathbf{A}x(0) - col_N(\bar{x})}{\mathbf{Q}}
		} \green{ + 
			\Big[
				\mathbf{A}x(0) - col_N(\bar{x})
			\Big]^T \mathbf{Q} \mathbf{Bu}
		} \\
		&\qquad \green{ + 
			\mathbf{u}^T \mathbf{B}^T \mathbf{Q}
			\Big[
				\mathbf{A}x(0) - col_N(\bar{x})
			\Big] 
		} \green{ + 
			\mathbf{u}^T \mathbf{B}^T \mathbf{Q} \mathbf{Bu}
		}
	\end{split}
\end{align}

\noindent
\begin{align}
	\blue{\bigcirc} &= \blue{
		\norma{\mathbf{u} - col_N(\bar{u})}{\mathbf{R}}
	} \\
	&= \blue{
		\Big[
			\mathbf{u} - col_N(\bar{u})
		\Big]^T \mathbf{R}
		\Big[
			\mathbf{u} - col_N(\bar{u})
		\Big]
	} \\
	&= \blue{
		\mathbf{u}^T \mathbf{R} \mathbf{u} 
		- \mathbf{u}^T \mathbf{R} \Big[ col_N(\bar{u}) \Big]
		- \Big[ col_N(\bar{u}) \Big]^T \mathbf{Ru}
		+ \norma{col_N(\bar{u})}{\mathbf{R}}
	}
\end{align}

substituindo $\green{\square}$ and $\blue{\bigcirc}$ de volta na equacao de custo e reorganizando os termos

\noindent
\begin{align}
	J(\mathbf{x}, \mathbf{u}) 
	&= \purple{
		\norma{x(0) - \bar{x}}{Q} 
		+ \norma{\mathbf{A} x(0) - col_N(\bar{x})}{\mathbf{Q}}
		+ \norma{col_N(\bar{u})}{\mathbf{R}}
	} \\
	&\quad + \orange{
		\Big[ \mathbf{A} x(0) - col_N(\bar{x}) \Big]^T \mathbf{QBu}
		+ \mathbf{u}^T \mathbf{B}^T \mathbf{Q} \Big[ \mathbf{A} x(0) - col_N(\bar{x}) \Big]
	} \\
	&\quad 
		- \mathbf{u}^T \mathbf{R} \Big[ col_N(\bar{u}) \Big] - \Big[ col_N(\bar{u}) \Big]^T \mathbf{Ru}
		+ \mathbf{u}^T (\mathbf{R} + \mathbf{B}^T \mathbf{Q} \mathbf{B}) \mathbf{u}
\end{align}

O \purple{termo constante} ser� denotado por \purple{$c_0$}

Lembrando que a matriz $\mathbf{Q}$ � sim�trica, tem-se

\begin{equation}
	\orange{ \left[\mathbf{A}x(0) - col_N(\bar{x})\right]^T \mathbf{QBu} } 
	= 
	\Big\{
	\orange{ \mathbf{u}^T \mathbf{B}^T \mathbf{Q} \left[\mathbf{A}x(0) - col_N(\bar{x})\right] }
	\Big\}^T \\
\end{equation}

Como esses termos s�o escalares, conclui-se que s�o iguais entre si. Logo:

\begin{equation}
	\begin{gathered}
		\orange{ \left[\mathbf{A}x(0) - col_N(\bar{x})\right]^T \mathbf{QBu} }
		+ \orange{ \mathbf{u}^T \mathbf{B}^T \mathbf{Q} \left[\mathbf{A}x(0) - col_N(\bar{x})\right] } \\
		= \orange{ 2 \left[\mathbf{A}x(0) - col_N(\bar{x})\right]^T \mathbf{QBu} }
	\end{gathered}
\end{equation}

De forma similar:

\noindent
\begin{equation}
	- \mathbf{u}^T \mathbf{R} \left[col_N(\bar{u})\right] - \left[col_N(\bar{u})\right]^T \mathbf{Ru} = 
	-2 \left[col_N(\bar{u})\right]^T \mathbf{Ru}
\end{equation}

Finalmente, pode-se escrever

\noindent
\begin{equation}
	\begin{gathered}
		J(\mathbf{u}) = \mathbf{u}^T (\mathbf{B}^T \mathbf{Q} \mathbf{B} + \mathbf{R}) \mathbf{u} \\
		+ 2 \left\{
			\left[
				\mathbf{A}x(0) - col_N(\bar{x})
			\right]^T \mathbf{QB} - \left[col_N(\bar{u})\right]^T \mathbf{R}
		\right\} \mathbf{u} = c_0
	\end{gathered}
\end{equation}

com 

\noindent
\begin{equation}
	c_0 = \norma{x(0) - \bar{x}}{Q} + \norma{\mathbf{A}x(0) - col_N(\bar{x})}{\mathbf{R}} + \norma{col_N(\bar{u})}{\mathbf{R}}
\end{equation}

Note-se que passaremos a usar a nota��o $J(\mathbf{u})$ em lugar de $J(\mathbf{x}, \mathbf{u})$, uma vez que $\mathbf{x}$ foi eliminada da express�o do custo.

Para terminar a reformula��o do problema de otimiza��o em termos de $\mathbf{u}$, deve-se considerar o v�nculo entre $\mathbf{x}$ e $\mathbf{u}$ nas \textbf{restri��es} 
